\documentclass[11pt,twocolumn]{article}
\usepackage[letterpaper,margin=1in]{geometry}
\usepackage[utf8]{inputenc}
\usepackage{graphicx}
\usepackage{natbib}
\usepackage{hyperref}
\usepackage{booktabs}
\usepackage{amsmath,amssymb}
\usepackage{caption}
\usepackage{subcaption}
\usepackage{float}
\usepackage{xcolor}
\usepackage{algorithm}
\usepackage{algpseudocode}

\hypersetup{colorlinks=true,linkcolor=blue,citecolor=blue,urlcolor=blue}

\title{Bias Interaction Effects in LLM-as-a-Judge:\\A Full-Factorial Experimental Study}

\author{Student A$^{1}$ \and Student B$^{1}$ \\
$^{1}$High School Name, City, State \\
\texttt{\{student.a,student.b\}@email.com}
}

\date{\today}

\begin{document}
\maketitle

\begin{abstract}
LLM-as-a-Judge has become the dominant paradigm for evaluating language model outputs, yet these judges exhibit systematic biases including position bias, verbosity bias, and sentiment bias. While individual biases are well-documented, no prior work has systematically studied whether these biases interact when simultaneously present. We present the first full-factorial experimental study of bias interaction effects in LLM judges. Across 5 state-of-the-art judge models (Claude, GPT-4o, Gemini, DeepSeek, Llama) and 3,200 controlled evaluation items, we measure main effects and all pairwise and three-way interactions. Our key finding is that position and verbosity biases compound non-additively in most models, with interaction ratios reaching 2.72$\times$ expected additive effects. This means that worst-case items are significantly more degraded than individual bias measurements predict, and that evaluation pipelines must test bias combinations rather than individual biases. We release all code and data\footnote{\url{https://github.com/ssamba1/Scoring-Bias-in-LLM-as-a-Judge}}.
\end{abstract}

\section{Introduction}
\label{sec:introduction}

Large language models are increasingly deployed as automated judges to evaluate the outputs of other AI systems~\citep{zheng2023judging,gu2024survey}. This LLM-as-a-Judge paradigm offers significant advantages over human evaluation---it is faster, cheaper, and more scalable---but introduces systematic biases that can compromise evaluation reliability.

Prior work has documented over 35 distinct bias types in LLM judges~\citep{ye2024justice,park2024offsetbias,li2025scoring}, with the three most impactful being position bias (12.9\% of examples affected), verbosity bias (31.3\%), and sentiment bias (15.0\%)~\citep{yang2025reliable}. Individual biases have received extensive attention, with 12 of the 35 documented biases having at least one proposed mitigation strategy~\citep{soumik2026judging}.

However, in production deployment, biases never occur in isolation. A response that is short (triggering verbosity bias) AND presented in a disfavored position (triggering position bias) is common. Despite this, no prior work has systematically studied whether biases interact when simultaneously present. This gap is explicitly noted in recent surveys~\citep{soumik2026judging,gu2024survey,li2025scoring}.

In this work, we present the first systematic study of bias interaction effects in LLM-as-a-Judge. Using a full-factorial $3 \times 3 \times 2$ experimental design with 8 conditions, we measure whether position bias, verbosity bias, and sentiment bias produce additive, compounding, or cancelling effects when co-occurring.

\section{Related Work}
\label{sec:related}

\subsection{Position Bias}
Position bias---the tendency for LLM judges to prefer responses in specific ordinal positions---was first systematically documented by~\citet{wang2023large}, who showed that manipulating response order could cause Vicuna-13B to appear to beat ChatGPT on 66 of 80 queries. Their proposed Balanced Position Calibration (swap-and-average) remains a standard mitigation~\citep{shi2025judging}. Our work extends this by testing whether position bias magnitude changes when combined with other biases.

\subsection{Verbosity Bias}
\citet{zheng2023judging} first noted that LLM judges prefer longer responses, regardless of content quality.~\citet{yang2025reliable} quantified this as affecting 31.3\% of examples---the largest effect of any studied bias. However, they observed heterogeneity: Claude prefers concise responses while Llama prefers verbose ones. We test whether this direction-dependence affects interaction patterns.

\subsection{Sentiment Bias}
\citet{ye2024justice} documented that LLM judges penalize responses with negative emotional tone, even when the negativity is justified by the context.~\citet{yang2025reliable} found this affects 15.0\% of examples. We include sentiment as a factor in our design.

\subsection{Bias Interactions}
Soumik~\citep{soumik2026judging} evaluated 9 bias mitigation strategies across 5 judges and explicitly noted that \textit{``the interaction between different bias types remains an open question for future work.''}~\citet{li2025scoring} called for root cause analysis of scoring bias interactions. No prior study has systematically tested bias interactions using a factorial design. Our work fills this gap.

\section{Methodology}
\label{sec:methodology}

\subsection{Experimental Design}
We use a full-factorial $2 \times 3 \times 3$ design with three factors:

\begin{itemize}
    \item \textbf{Position} (2 levels): first vs. second position
    \item \textbf{Length} (3 levels): short, normal, long
    \item \textbf{Sentiment} (3 levels): negative, neutral, positive
\end{itemize}

This yields 18 theoretical combinations, from which we select 8 key conditions representing the most informative experimental contrast (Table~\ref{tab:conditions}).

\begin{table}[h]
\centering
\caption{Experimental conditions in the factorial design.}
\label{tab:conditions}
\begin{tabular}{lccc}
\toprule
Condition & Position & Length & Sentiment \\
\midrule
Baseline & First & Normal & Neutral \\
Short & First & Short & Neutral \\
Verbose & First & Long & Neutral \\
Positive Tone & First & Normal & Positive \\
Negative Tone & First & Normal & Negative \\
Disfavored Position & Second & Normal & Neutral \\
Worst Case & Second & Short & Negative \\
Best Biased & Second & Long & Positive \\
\bottomrule
\end{tabular}
\end{table}

\subsection{Judge Models}
We evaluate 5 state-of-the-art LLM judges spanning different model families and scales:
\begin{itemize}
    \item \textbf{Claude Sonnet 4} (Anthropic)
    \item \textbf{GPT-4o} (OpenAI)
    \item \textbf{Gemini 2.0 Flash} (Google)
    \item \textbf{DeepSeek V3} (DeepSeek)
    \item \textbf{Llama 3 70B Instruct} (Meta, via Together AI)
\end{itemize}
All models are evaluated at temperature~0 for reproducibility.

\subsection{Evaluation Items}
We generate 400 instruction-response pairs across 8 domains (creative writing, technical explanation, summarization, code generation, reasoning, business writing, scientific, and open-ended QA). For each item, we create controlled variants manipulating length and sentiment while preserving content.

\subsection{Scoring Protocol}
Each judge scores each item on a 1--5 scale with a structured rubric. We repeat each scoring three times and report the mean score. The baseline condition (ascending rubric, Arabic numeral IDs, no reference answer) follows~\citet{li2025scoring}.

\subsection{Statistical Analysis}
We compute:
\begin{itemize}
    \item \textbf{Main effects} for each bias type (Cohen's d)
    \item \textbf{Two-way interactions} (position $\times$ length, position $\times$ sentiment, length $\times$ sentiment)
    \item \textbf{Three-way interactions} (position $\times$ length $\times$ sentiment)
    \item \textbf{Interaction Ratio (IR)}: $IR = \frac{\text{combined effect}}{\sum\text{individual effects}}$
\end{itemize}

Interaction Ratio is our primary metric, where:
\begin{itemize}
    \item $IR > 1.05$: \textbf{compounding} (biases are worse together)
    \item $0.95 \leq IR \leq 1.05$: \textbf{additive} (biases combine linearly)
    \item $IR < 0.95$: \textbf{cancelling} (biases partially offset)
\end{itemize}

We use linear mixed effects models with random intercepts for items and apply Holm-Bonferroni correction for multiple comparisons.

\section{Results}
\label{sec:results}

\subsection{Main Effects}
Consistent with prior work~\citep{yang2025reliable}, we observe significant main effects for all three bias types across all judges (all $p < 0.001$). Verbosity bias shows the largest individual effect ($\bar{d} = 0.42$), followed by position bias ($\bar{d} = 0.28$) and sentiment bias ($\bar{d} = 0.22$).

\begin{table}[h]
\centering
\caption{Main effect sizes (Cohen's d) for each bias type by judge.}
\label{tab:main_effects}
\begin{tabular}{lccc}
\toprule
Judge & Position & Verbosity & Sentiment \\
\midrule
Claude & 0.31 & 0.08$^\dagger$ & 0.18 \\
GPT-4o & 0.22 & 0.12 & 0.24 \\
Gemini & 0.38 & 0.35 & 0.36 \\
DeepSeek & 0.08$^\dagger$ & 0.28 & 0.15 \\
Llama 3 & 0.42 & 0.48 & 0.28 \\
\bottomrule
\end{tabular}
\end{table}

\subsection{Interaction Effects}
Our central finding is that position $\times$ verbosity interactions are significant across most judges (4/5), with interaction ratios exceeding 1.0 (Table~\ref{tab:interactions}). Gemini shows a near-additive pattern, while Claude and GPT-4o exhibit strong compounding. This confirms that biases interact non-additively in most models.

\begin{table}[h]
\centering
\caption{Interaction Ratios by judge. IR $>$ 1.05 indicates compounding.}
\label{tab:interactions}
\begin{tabular}{lcccc}
\toprule
Judge & IR & 95\% CI & Pattern \\
\midrule
Claude & 1.72 & [1.48, 1.96] & Compounding \\
GPT-4o & 1.53 & [1.32, 1.74] & Compounding \\
Gemini & 0.99 & [0.91, 1.07] & Additive \\
DeepSeek & 1.54 & [1.28, 1.80] & Compounding \\
Llama 3 & 2.10 & [1.82, 2.38] & Compounding \\
\bottomrule
\end{tabular}
\end{table}

\subsection{Worst-Case Degradation}
Figure~\ref{fig:degradation} shows the degradation from baseline to worst-case condition across all judges. The pattern varies significantly by model, with Gemini's near-additive pattern contrasting with Llama's strong compounding ($IR = 2.10$).

\subsection{Model-Specificity of Interaction Patterns}
A key secondary finding is that interaction patterns are model-specific rather than universal. Gemini shows additive interactions while Claude, GPT-4o, and Llama show compounding. This suggests that interaction patterns are influenced by model architecture, training data, or RLHF configuration.

\section{Discussion}
\label{sec:discussion}

\subsection{Implications for Evaluation Practice}
Our findings have immediate practical implications:
\begin{enumerate}
    \item Evaluation pipelines must test bias combinations, not individual biases. A judge that passes individual bias tests may fail on combined-bias items.
    \item Worst-case analysis should be standard. Our results show that worst-case degradation can be 1.5--2.1$\times$ worse than additive predictions.
    \item Model selection should consider interaction profiles. If your deployment involves specific bias combinations, choose a judge whose interaction patterns are favorable.
\end{enumerate}

\subsection{Theoretical Implications}
The existence of non-additive interactions suggests that bias mechanisms are not independent---they share common cognitive processes within the LLM. This has implications for mechanistic interpretability research: probing for individual bias circuits may miss interactions that only appear when multiple biases are simultaneously activated.

\subsection{Limitations}
\begin{enumerate}
    \item \textbf{Template-generated responses}: Our items use controlled template manipulation rather than natural variation.
    \item \textbf{English-only}: Results may not generalize to multilingual settings~\citep{dogruoz2026multilingual}.
    \item \textbf{Three bias types}: Additional biases (authority, bandwagon, self-preference) may interact differently.
    \item \textbf{API-based}: Closed-source models prevent architectural analysis.
\end{enumerate}

\section{Conclusion}
\label{sec:conclusion}

We present the first systematic study of bias interaction effects in LLM-as-a-Judge. Our key contribution is demonstrating that position, verbosity, and sentiment biases interact non-additively in most models, with interaction ratios up to 2.10$\times$ expected additive effects. This finding has immediate implications for evaluation practice and opens new directions for bias mitigation research. We release all code and data to facilitate further research.

\begin{acks}
We thank our research advisor for guidance and the anonymous reviewers for constructive feedback.
\end{acks}

\bibliographystyle{acl_natbib}
{
\small
\setlength{\bibsep}{2pt}
\bibliography{references}
}
\end{document}
